\documentclass[11pt]{article}
\usepackage[utf8]{inputenc}
\usepackage[spanish]{babel}
\usepackage{graphicx}         % Para incluir imágenes
\usepackage{amsmath}          % Para notación matemática
\usepackage[margin=3cm]{geometry}  % Márgenes
\usepackage[font=normalsize,labelfont=small]{caption}  % Estilo de captions
\usepackage{hyperref}         % Para links clickeables

%%%%%%%%%%%%%%%%%%%%%%%%%%%%%%%%%%%%%%%%%%%%%
% EXTENSIÓN MÁXIMA: 10 HOJAS SIN EXCEPCIÓN. %
%%%%%%%%%%%%%%%%%%%%%%%%%%%%%%%%%%%%%%%%%%%%%

\begin{document}

\begin{titlepage}
    \centering
    \vspace*{2cm}
    \includegraphics[scale=1.8]{Logo-Udesa.jpg}\par
    \vspace{10pt}

    {\LARGE \textbf{I302 - Aprendizaje Automático\\ y Aprendizaje Profundo}\par}
    \vspace{1cm}

    {\LARGE \textbf{Trabajo Práctico 4}\par}
    \vspace{4cm}
    
    {\LARGE {Abril Diaco}\par}  % <- Modificar Nombre y Apellido
    \vspace{4cm}
    
    {\Large \today\par}
    \vspace{1cm}
    \Large{Ingeniería en Inteligencia Artificial}
\end{titlepage}

% RESUMEN
\begin{abstract}
Breve resumen (unas líneas) de qué se hizo, cómo se hizo, qué resultados se obtuvieron. Por ejemplo: ``En este trabajo se quizo modelar... Para llevarlo a cabo, se desarrolló ... se probó utilizando ... . El modelo performó ... y las métricas dieron ...''.
\end{abstract}

% INTRODUCCION
\section{Introducción}
\textit{Explicar el problema. Proporcionar el contexto necesario y explicar cuáles son las entradas y salidas del algoritmo. Por ejemplo: “La entrada del algoritmo es (imagen, amplitud de señal, edad del paciente, mediciones de lluvia, video en escala de grises, etc.). Luego, se utilizó un/a (red neuronal, regresión lineal, árbol de decisión, etc.) para obtener una predicción de (edad, precio de acciones, tipo de cáncer, género musical, etc.)”.}

% METODOS
\section{Métodos}
\textit{Describir los algoritmos que se implementaron y utilizaron. Asegurarse de incluir la notación matemática relevante. Si el espacio lo permite, se podrá dar una breve descripción ($\approx$ 1 párrafo) de cómo funciona.}

\textit{Describir el conjunto de datos: ¿Qué proporción de entrenamiento/validación/prueba se utilizó? ¿Se realizó algún preprocesamiento? ¿Qué tipo de normalización o aumento de datos se utilizó?}
% %
% El conjunto de datos se encuentra perfectamente balanceado: cada una de las diez clases contiene 2500 observaciones, representando el 10\% del total. Esto resulta conveniente para el análisis posterior, ya que evita que los algoritmos de agrupamiento se vean influenciados por diferencias en la frecuencia de las clases.

% Al observar muestras agrupadas por clase, se aprecia una variabilidad considerable dentro de una misma categoría. Por ejemplo, las clases de calzado presentan diferencias en orientación, altura y forma, mientras que la clase de bolsos incluye objetos con diseños y tamaños diversos. Esta variabilidad intraclase sugiere que la estructura de los datos puede no coincidir perfectamente con las etiquetas originales.

% Las imágenes promedio por clase permiten identificar la silueta característica de cada categoría. Algunas clases, como pantalón, zapatilla, bolso y bota, presentan formas promedio bien diferenciadas. En cambio, clases como remera/top, pulóver, abrigo y camisa resultan visualmente más similares entre sí, lo cual podría dificultar su separación mediante métodos no supervisados.

% Finalmente, las estadísticas de intensidad muestran que las clases no ocupan la imagen de la misma manera. Las sandalias y zapatillas presentan menor intensidad promedio y menor cantidad de píxeles no nulos, mientras que prendas como pulóver, abrigo y camisa ocupan una mayor proporción del área de la imagen. Estas diferencias pueden influir en los agrupamientos, ya que los modelos pueden capturar tanto la forma de los objetos como la cantidad de píxeles activos.
% %
% se uso un random stratificado con 20\% de los datos para evaluación y el resto para entrenamiento.
%Se utilizó una semilla fija para garantizar la reproducibilidad de los resultados.

\textit{Detallar los hiperparámetros elegidos y cómo se seleccionaron. ¿Realizaron validación cruzada? Si es así, ¿cuántos folds utilizaron?}

\textit{Incluir métricas de rendimiento utilizadas.}

% RESULTADOS
\section{Resultados}
\textit{Mostrar y discutir los resultados más destacables del trabajo. ¿Qué modelo funcionó mejor? ¿Qué modelo funcionó peor? Discutir porqué cree que fue así.}

\textit{Utilice gráficos relevantes para mostrar los puntos clave de sus resultados. Estos deben llevar leyenda, labels, etc., adecuados para que el lector entienda estos gráficos.}

\end{document}




